\mode*

\section{Disposition}

\begin{frame}
  \begin{figure}
    \includegraphics[height=0.8\textheight]{fig/book.jpg}
    \caption{Front cover of Necessary Conditions of Learning.}
  \end{figure}
\end{frame}

\begin{frame}
  \begin{enumerate}
    \item What makes humans human?
    \alert{\item What is to be learned?}
    \item Sameness and difference in learning
    \item What does the world look like to others?
    \item The art of learning
    \item Making learning possible
    \item Learning to help others to learn
  \end{enumerate}
\end{frame}

\begin{frame}
  \begin{center}
    What is to be learned?\\[1em]
    instead of\\[1em]
    What is to be done?
  \end{center}

  \begin{block}{Ch 2 What is to be learned?}
    \begin{itemize}
      \item What is an intended learning outcome?
      \item What should it be?
      \item It depends on the learner: \emph{critical aspects, critical 
        features} --- the atoms of the learning objective.
    \end{itemize}
  \end{block}
  
  \pause

  \begin{block}{The goal}
    \begin{itemize}
      \item \textcquote{NecessaryConditionsOfLearning}{The focus of the theory 
        elaborated in this book is on learning to handle new situations in 
      powerful ways.}
    \end{itemize}
  \end{block}
\end{frame}

\section{Critial aspects and features}

\begin{frame}
  \begin{figure}
    \begin{subfigure}{0.3\columnwidth}
      \centering
      \includegraphics{fig/contrast-color.tikz}
      \caption{Contrast}
    \end{subfigure}
    \hfill
    \begin{subfigure}{0.3\columnwidth}
      \centering
      \includegraphics{fig/generalization-color.tikz}
      \caption{Generalization}
    \end{subfigure}
    \hfill
    \begin{subfigure}{0.3\columnwidth}
      \centering
      \includegraphics{fig/fusion-color.tikz}
      \caption{Fusion}
    \end{subfigure}
    \caption{%
      Illustrating the patterns of variation for aspects color and shape.
    }
  \end{figure}

  \begin{onlyenv}<1>
    \begin{example}
      \begin{itemize}
        \item Critical aspect: colour.
        \item Critical feature: blue.
        \item Non-critical feature: green.
        \item Non-critical aspect: shape.
        \item Non-critical features: circle, square.
      \end{itemize}
    \end{example}
  \end{onlyenv}
  \begin{onlyenv}<2>
    \begin{remark}
      \begin{itemize}
        \item There are other aspects too, \emph{unintentionally}.
        \item Another aspect in this figure is order: the circle is always 
          \enquote{first} (to the right).
      \end{itemize}
    \end{remark}
  \end{onlyenv}
\end{frame}

\begin{frame}
  \begin{figure}
    \includegraphics[height=0.5\textheight]{fig/maja-writing.png}
    \caption{My daughter's writing. (Corresponds to~\cite[Fig.~2.1, 
    p.~30]{NecessaryConditionsOfLearning}.)}
  \end{figure}
  \begin{example}
    \begin{itemize}
      \item Obviously my daughter knows some aspect(s) of writing.
      \item Just not all that are needed.
    \end{itemize}
  \end{example}
\end{frame}

\begin{frame}
  \blockcquote[p.~37]{NecessaryConditionsOfLearning}
  {%
    The object of learning \textelp{} amounts to becoming able to discern all 
    the critical aspects and to focus on them simultaneously.%
  }
\end{frame}

